\documentclass[11pt]{article}
\usepackage[letterpaper,margin=0.65in,includehead,includefoot,headheight=15pt]{geometry}
\usepackage{amsmath,amssymb}
\usepackage{enumitem}
\usepackage{fancyhdr}
\usepackage{microtype}
\usepackage[hidelinks]{hyperref}
\setlength{\parindent}{0pt}
\setlength{\parskip}{4pt}
\setlist[enumerate]{leftmargin=*,itemsep=3pt,topsep=2pt}
\pagestyle{fancy}
\fancyhf{}
\fancyhead[L]{\textbf{Induction Core}}
\fancyhead[R]{COMP 2300}
\fancyfoot[C]{\thepage}
\renewcommand{\headrulewidth}{0.4pt}
\fancypagestyle{plain}{\fancyhf{}\fancyhead[L]{\textbf{Induction Core}}\fancyhead[R]{COMP 2300}\fancyfoot[C]{\thepage}\renewcommand{\headrulewidth}{0.4pt}}

\begin{document}
\textbf{Name:} \rule{3.2in}{0.4pt}\hfill
\textbf{Attempt date:} \rule{1.55in}{0.4pt}

\textbf{Time: 20 minutes.} A supervised form contains five labeled problems,
one from each section A--E. Each label is one problem, including all of its
parts. \emph{Meets} requires four fully correct problems; \emph{Exceeds}
requires all five. Every retake uses a new complete form. The problems below
are practice variants; a live form contains only one problem from each section.

\textbf{Directions.} Use the requested proof method. State the proposition and
the relevant hypothesis precisely, verify every required base case, and identify
the exact line where the hypothesis advances the proof.

\section*{A. Proof Architecture}

\subsection*{A.1 Summation Blueprint}
For the claim $1+2+\cdots+n=n(n+1)/2$ for $n\ge1$, write: (i) $P(1)$,
(ii) the induction hypothesis $P(k)$, and (iii) the exact statement $P(k+1)$
that must be proved. Do not complete the algebra.

\subsection*{A.2 Divisibility Blueprint}
For the claim $3\mid(n^3+2n)$ for every $n\ge1$, write the base case, the
induction hypothesis using an integer witness, and the divisibility statement
needed for the next case. Do not complete the algebra.

\subsection*{A.3 Strong-Induction Blueprint}
For the claim that every integer $n\ge2$ is a product of primes, state the base
case, the strong induction hypothesis for an arbitrary $k\ge2$, and what must
be proved for $k+1$. Explain in one sentence why ordinary induction is less
natural here.

\subsection*{A.4 Structural-Induction Blueprint}
A full binary tree is either one leaf or a root joined to two smaller full
binary trees. To prove that every full binary tree has one more leaf than
internal vertex, state the basis case, the two structural induction hypotheses,
and the statement required for the newly joined tree. Do not complete the count.

\subsection*{A.5 Choose the Method}
Choose ordinary, strong, or structural induction for each claim and give one
sentence describing the hypothesis.
\begin{enumerate}[label=(\alph*)]
  \item $1+3+\cdots+(2n-1)=n^2$.
  \item Every amount $n\ge18$ can be made using 4-cent and 7-cent stamps.
  \item Every recursively formed parenthesized expression has equally many
        left and right parentheses.
\end{enumerate}

\newpage
\section*{B. Ordinary Induction}

\subsection*{B.1 Odd-Number Sum}
Prove by ordinary induction that
$1+3+\cdots+(2n-1)=n^2$ for every $n\ge1$.

\subsection*{B.2 Exponential Divisibility}
Prove by ordinary induction that $7\mid(8^n-1)$ for every $n\ge0$.

\subsection*{B.3 Triangular Numbers}
Prove by ordinary induction that
$1+2+\cdots+n=n(n+1)/2$ for every $n\ge1$.

\subsection*{B.4 Consecutive Product}
Prove by ordinary induction that $2\mid(n^2+n)$ for every $n\ge1$.

\subsection*{B.5 Geometric Sum}
Prove by ordinary induction that
\[
  1+3+3^2+\cdots+3^n=\frac{3^{n+1}-1}{2}
\]
for every $n\ge0$.

\newpage
\section*{C. Strong Induction}

\subsection*{C.1 Postage}
Using the verified base cases $18,19,20,21$, complete the strong-induction step
showing that every amount $n\ge18$ can be made from 4-cent and 7-cent stamps.

\subsection*{C.2 Prime Products}
Prove by strong induction that every integer $n\ge2$ is a product of one or
more primes. In the inductive step, separate the cases where $k+1$ is prime or
composite.

\subsection*{C.3 Distinct Powers of Two}
Assume every integer from $1$ through $k$ is a sum of distinct nonnegative
powers of $2$. Let $2^m$ be the largest power of $2$ not exceeding $k+1$.
Complete the step proving the claim for $k+1$, including the zero-remainder case.

\subsection*{C.4 Jigsaw Moves}
A jigsaw puzzle is assembled by joining two existing blocks at each move. Use
strong induction to prove that assembling $n\ge1$ pieces always takes exactly
$n-1$ moves, regardless of the order of joins.

\subsection*{C.5 Chocolate Breaks}
A rectangular chocolate bar has $n$ unit squares. One move breaks one existing
rectangular piece into two smaller rectangular pieces. Use strong induction to
prove that separating all $n$ squares always takes exactly $n-1$ breaks.

\newpage
\section*{D. Recursive Objects}

\subsection*{D.1 Linear Recurrence}
Let $a_0=2$ and $a_{n+1}=3a_n+2$. Compute $a_1,a_2,a_3$, conjecture a closed
formula, and verify only the induction step for that formula.

\subsection*{D.2 Arithmetic Sequence}
Give a recursive definition of $a_n=4n-2$ for $n\ge1$. Then state the
induction hypothesis and verify the next-case algebra.

\subsection*{D.3 Parenthesized Expressions}
Variables belong to $\mathcal E$; if $u,v\in\mathcal E$, then $(u+v)$ and
$(u\times v)$ belong to $\mathcal E$. Prove by structural induction that
every expression in $\mathcal E$ has equally many left and right parentheses.

\subsection*{D.4 Full Binary Trees}
A full binary tree is one leaf or a root joined to two full binary trees. Prove
by structural induction that every such tree has one more leaf than internal
vertex.

\subsection*{D.5 String Reversal}
Let $\lambda^R=\lambda$ and $(wx)^R=xw^R$ for a string $w$ and symbol $x$.
Using $\ell(\lambda)=0$ and $\ell(wx)=\ell(w)+1$, prove by structural
induction that $\ell(w^R)=\ell(w)$ for every string $w$.

\newpage
\section*{E. Diagnosis and Synthesis}

\subsection*{E.1 Circular Step}
A student claims to prove $2^n\ge n+1$ for $n\ge0$ by assuming
$2^{k+1}\ge k+2$ and declaring the next case complete. Identify the error,
give the correct base case and induction hypothesis, and show the valid
next-case inequality.

\subsection*{E.2 Missing Base Case}
A recurrence proof uses both $P(k)$ and $P(k-1)$ to establish $P(k+1)$ but
checks only $P(0)$. Explain precisely what is missing and state a sufficient
pair of base cases and strong induction hypothesis.

\subsection*{E.3 A False Induction Proof}
A student claims all positive integers are equal. The base case with one integer
is true. In the step, the student applies the hypothesis to the first $k$ and
last $k$ integers in a list of $k+1$ integers, then says the overlapping entries
connect the two groups. Identify the first value of $k$ for which this overlap
argument is needed and explain why the step fails there.

\subsection*{E.4 Binomial-Sum Step}
Assume $\sum_{j=0}^{k}\binom{k}{j}=2^k$. Use Pascal's identity
$\binom{k+1}{j}=\binom{k}{j}+\binom{k}{j-1}$ to prove
$\sum_{j=0}^{k+1}\binom{k+1}{j}=2^{k+1}$. Show how the two sums are reindexed.

\subsection*{E.5 Repair Distinct-Powers Reasoning}
A student tries to prove that every positive integer is a sum of distinct
powers of $2$ by representing $k$ and then adding $1$. Give a concrete value
of $k$ for which this repeats a power, and replace the step with a valid strong-
induction argument using the largest power of $2$ not exceeding $k+1$.

\end{document}
